第二章的两次 POST 分别提交 P1、P2，各自保住了订单和流水关系。两次 HTTP 尝试之间却缺少一项关系：服务没有收到能够说明它们属于同一次购买的信息。商品、账户和金额相同，仍不足以补出这项关系；Alice 也可能有意连续买两份相同商品。

本章继续用 HTTP 报文表示购买输入，逐步修改 API 的身份规则：请求要带什么标识、重试怎样沿用它，以及服务如何据此区分恢复与新购买。

\textbf{这里仍然只处理一个商品、数量为一、以 CNY 分计价的一次购买。}客户端先为用户确认的购买保留一个身份，以后的查询和重试继续携带它。服务由此可以查找原购买，核对输入，再决定返回原结果还是受理新购买。身份需要在什么时候产生、怎样保存、在哪个范围内解释，是设计这套处理规则时首先遇到的几个问题。
